\chapter*{Introduction Générale}
\addcontentsline{toc}{chapter}{INTRODUCTION GÉNÉRALE}
\adjustmtc
\thispagestyle{MyStyle}

Dans un environnement économique marqué par une forte concurrence et par une exigence croissante en matière de qualité de service, la gestion de la relation client s'impose comme un enjeu stratégique majeur. Les entreprises ne peuvent plus se contenter d'outils dispersés pour suivre leurs prospects, leurs opportunités commerciales, leurs interactions de support et leurs engagements contractuels. Elles ont besoin d'un système centralisé, capable de structurer l'information, d'améliorer la coordination entre équipes et de fournir une vision fiable de l'activité en temps réel.\par
\vspace{0.3cm}

Les solutions CRM du marché répondent en partie à cette problématique, mais elles restent souvent coûteuses, surdimensionnées ou difficiles à adapter aux contraintes spécifiques d'une entreprise donnée. Dans ce contexte, le projet réalisé au sein de \textbf{FRS} vise à concevoir une solution maîtrisée, personnalisable et cohérente avec les besoins réels de l'organisation. L'ambition n'est pas seulement de reproduire des fonctionnalités classiques de prospection ou de ticketing, mais de bâtir un écosystème numérique intégrant à la fois le travail interne, le suivi client et l'accès à l'information métier.\par
\vspace{0.3cm}

Le travail mené durant ce stage aboutit ainsi à la réalisation d'un système CRM articulé autour de quatre composantes complémentaires :
\begin{itemize}
    \item un \textbf{CRM web interne} pour les équipes administratives, commerciales, support et managériales ;
    \item un \textbf{portail client web} dédié à la consultation des tickets, devis et informations personnelles ;
    \item une \textbf{application mobile} orientée portail client pour l'accès nomade aux fonctionnalités essentielles ;
    \item un \textbf{assistant conversationnel} fondé sur une base documentaire et la technique RAG.
\end{itemize}

L'application web couvre les principaux besoins métier d'un CRM moderne : authentification et gestion des rôles, suivi des leads, création des comptes et contacts, pilotage du pipeline commercial, gestion des devis, traitement des tickets, suivi des activités, journalisation des actions et génération de tableaux de bord. Le portail client et l'application mobile prolongent cette logique côté utilisateur final, tandis que le chatbot apporte une couche d'assistance documentaire pour fluidifier l'accès à l'information.\par
\vspace{0.3cm}

La réalisation du projet s'inscrit dans une démarche \textbf{Agile Scrum} structurée en quatre sprints successifs :
\begin{itemize}
    \item \textbf{Sprint 01} : étude préalable, analyse des besoins, identification des acteurs et élaboration du backlog ;
    \item \textbf{Sprint 02} : conception détaillée, modélisation UML, architecture technique et conception de l'API ;
    \item \textbf{Sprint 03} : réalisation technique du backend, du frontend, de l'application mobile et du module d'assistance ;
    \item \textbf{Sprint 04} : tests, validation, corrections de stabilisation et préparation du déploiement local conteneurisé.
\end{itemize}

Le présent rapport est organisé en cinq chapitres. Le premier chapitre présente le cadre général du projet et son environnement. Le deuxième chapitre expose l'étude préalable et l'analyse des besoins. Le troisième chapitre détaille la phase de conception. Le quatrième chapitre décrit la réalisation technique des différentes composantes. Enfin, le cinquième chapitre est consacré aux tests, à la validation et au déploiement.\par
